% !TEX root = ../presentation.tex
% =============================================================
%  Frame — the line as a two-port: the nominal pi-section
%  Part of ac-transmission; \input from ac-transmission.tex,
%  after symbols.tex -- it opens the chapter's argument.
%
%  First step: the line is a two-port, not one impedance.  Series R'
%  and L', shunt C'.  power-transfer.tex then solves it.
%  Figure fig/pi-section.tex (inline TikZ).  Worked values for this
%  line are collected on reference-line.tex.
% =============================================================
\providecommand{\acfig}[1]{ac-transmission/fig/#1}

\begin{frame}{The line as a two-port --- the $\pi$-section}
  \begin{block}{Where the capacitance enters}
    {The line also charges itself\\
    Series $R'$ and $L'$, shunt $C'$ --- a two-port, not one impedance.}
  \end{block}

  \vspace{0.5em}
  \begin{columns}[c,onlytextwidth]
    \begin{column}{0.44\textwidth}\centering
      \input{\acfig{pi-section.tex}}\\[0.8em]
      \setlength{\abovedisplayskip}{2pt}\setlength{\belowdisplayskip}{2pt}
      \[ \begin{aligned}
           \underline{Z} &= \left(R'+j\cdot\omega\cdot L'\right)\cdot l \\
           \underline{Y} &= \left(G'+j\cdot\omega\cdot C'\right)\cdot l
         \end{aligned} \]
    \end{column}
    \begin{column}{0.54\textwidth}
      \setlength{\abovedisplayskip}{2pt}\setlength{\belowdisplayskip}{2pt}
      \[ \begin{bmatrix}\underline{U}_S\\[0.2em] \underline{I}_S\end{bmatrix}
         = \begin{bmatrix}\underline{A} & \underline{B}\\[0.2em]
                          \underline{C} & \underline{D}\end{bmatrix}
           \cdot
           \begin{bmatrix}\underline{U}_R\\[0.2em] \underline{I}_R\end{bmatrix} \]
      \[ \begin{aligned}
           \underline{A}=\underline{D} &= 1+\frac{\underline{Z}\cdot\underline{Y}}{2} \\
           \underline{B} &= \underline{Z} \\
           \underline{C} &= \underline{Y}\cdot
                 \left(1+\frac{\underline{Z}\cdot\underline{Y}}{4}\right)
         \end{aligned} \]
    \end{column}
  \end{columns}

\end{frame}
